\documentclass[a4paper,12pt]{article}

% ------------------- PACKAGES -------------------
\usepackage[utf8]{inputenc}
\usepackage[T1]{fontenc}
\usepackage[english]{babel}
\usepackage{graphicx}
\usepackage{hyperref}
\usepackage{geometry}
\usepackage{fancyhdr}
\usepackage{setspace}
\usepackage{titlesec}
\usepackage{booktabs}
\usepackage{float}
\usepackage{xcolor}
\usepackage{listings}
\usepackage{amsmath}
\usepackage{array}

% ------------------- CODE LISTING SETTINGS -------------------
\definecolor{codegray}{rgb}{0.5,0.5,0.5}
\definecolor{codepurple}{rgb}{0.58,0,0.82}
\definecolor{backcolour}{rgb}{0.95,0.95,0.92}

\lstdefinestyle{mystyle}{
    backgroundcolor=\color{backcolour},   
    commentstyle=\color{codegray},
    keywordstyle=\color{magenta},
    numberstyle=\tiny\color{codegray},
    stringstyle=\color{codepurple},
    basicstyle=\ttfamily\footnotesize,
    breakatwhitespace=false,         
    breaklines=true,                 
    captionpos=b,                    
    keepspaces=true,                 
    numbers=left,                    
    numbersep=5pt,                  
    showspaces=false,                
    showstringspaces=false,
    showtabs=false,                  
    tabsize=2
}
\lstset{style=mystyle}

% ------------------- PAGE SETTINGS -------------------
\geometry{margin=2.5cm}
\setstretch{1.25}
\hypersetup{
    colorlinks=true,
    linkcolor=blue,
    urlcolor=blue,
    citecolor=blue
}
\pagestyle{fancy}
\fancyhf{}
\fancyhead[L]{Parallelization Benchmark – Task 3}
\fancyhead[R]{Big Data – ULPGC}
\fancyfoot[C]{\thepage}
\fancyfoot[R]{December 6, 2025}

% ------------------- TITLE FORMATTING -------------------
\titleformat{\section}{\large\bfseries}{\thesection.}{1em}{}
\titleformat{\subsection}{\normalsize\bfseries}{\thesubsection.}{1em}{}

% ------------------- DOCUMENT START -------------------
\begin{document}

% ------------------------------------------------------------
% Cover Page
% ------------------------------------------------------------
\begin{titlepage}
    \centering
    \includegraphics[width=0.3\textwidth]{ulpgc_logo.png}\par\vspace{1cm}
    {\Large\textbf{Universidad de Las Palmas de Gran Canaria}}\\[0.3cm]
    {\large Faculty of Computer Science and Engineering}\\[1cm]
    {\Huge\bfseries Task 3: Parallelization Benchmark of Matrix Multiplication\par}
    \vspace{0.8cm}
    {\Large Big Data – Academic Year 2025–2026\par}
    \vspace{1.5cm}
    {\Large \textbf{Individual Assignments}\par}
    \vspace{0.3cm}
    {\large \textbf{Author:}\par}
    \begin{center}
    \begin{tabular}{ll}
    Adrián Ojeda Viera & (DNI: 45364780V) \\
    \end{tabular}
    \end{center}
    \vspace{0.8cm}
    {\large \textbf{Repository:}\par}
    \vspace{0.2cm}
    \href{https://github.com/adriaanojeda/BIGDATA.git}{\textcolor{blue}{https://github.com/adriaanojeda/BIGDATA.git}}
    \vfill
    {\large \textbf{Submission Date:} December 6, 2025\par}
    \vspace{1cm}
    \rule{0.9\textwidth}{0.4pt}\\[0.3cm]
    {\small \textit{“This document analyzes the impact of parallel computing techniques (Streams, Threads, Executors) and memory-layout optimizations (Vectorization) on the performance of dense matrix multiplication.”}}
\end{titlepage}

% ------------------------------------------------------------
% Abstract
% ------------------------------------------------------------
\begin{abstract}
This report presents a comparative study of different matrix multiplication implementations in Java. We evaluated sequential, parallel (using Streams and manual Threads), and memory-optimized (simulated vectorization) approaches. Benchmarks conducted on matrix sizes up to $N=4096$ reveal that while parallelization offers significant speedups ($\sim 9x$), the single-threaded memory-optimized approach outperforms all parallel versions with a speedup of $\sim 38x$. This finding underscores that for memory-bound algorithms like matrix multiplication, data locality and cache usage are often more critical than raw core count.
\end{abstract}

\newpage

% ------------------------------------------------------------
\section{Introduction and Objectives}

Matrix multiplication is a computationally intensive operation with a complexity of $O(n^3)$. In the context of Big Data, optimizing this operation is crucial. Modern CPUs offer multiple cores and vector instructions (SIMD), providing opportunities for parallelization.

The main objectives of this task were:
\begin{itemize}
    \item \textbf{Implementation}: Develop parallel versions of matrix multiplication using Java's concurrency tools (Parallel Streams, Threads, Atomic variables).
    \item \textbf{Optimization}: Implement a "Vectorized" version that improves memory access patterns (Cache Friendly).
    \item \textbf{Benchmarking}: Measure Execution Time, Speedup, and Efficiency for large matrices ($N=4096$).
\end{itemize}

% ------------------------------------------------------------
\section{Methodology}

\subsection{Experimental Setup}
The experiments were executed on a machine with the following specifications:
\begin{itemize}
    \item \textbf{CPU:} AMD Ryzen 5 3600X 6-Core Processor % Rellena esto
    \item \textbf{Cores/Threads:} 6 Cores / 12 Threads % Rellena esto (ej: 4 Cores / 8 Threads)
    \item \textbf{RAM:} 32GB % Rellena esto
    \item \textbf{Java Version:} java 23.0.2 2025-01-21
\end{itemize}

\subsection{Implementations}

We compared five different approaches:

\begin{enumerate}
    \item \textbf{Sequential:} The standard triple-loop algorithm ($ijk$ order).
    \item \textbf{Parallel Streams:} Uses \texttt{IntStream.range().parallel()} to distribute the outer loop across the ForkJoinPool.
    \item \textbf{Threads (Manual):} Explicitly creates one \texttt{Thread} per row.
    \item \textbf{Atomic Bag of Tasks:} Uses a fixed thread pool where workers dynamically fetch row indices from an \texttt{AtomicInteger}.
    \item \textbf{Vectorized (Simulated):} A single-threaded version that transposes Matrix B before multiplication to allow contiguous memory access ($ikj$ access pattern equivalent).
\end{enumerate}

% ------------------------------------------------------------
\section{Implementation Details}

\subsection{Parallel Streams Approach}
Java 8 Streams provide a high-level abstraction for parallelism. This implementation is concise and delegates thread management to the JVM.

\begin{lstlisting}[language=Java, caption={Parallel Stream Implementation}]
IntStream.range(0, n).parallel().forEach(i -> {
    for (int j = 0; j < n; j++) {
        double sum = 0.0;
        for (int k = 0; k < n; k++) {
            sum += A[i][k] * B[k][j];
        }
        C[i][j] = sum;
    }
});
\end{lstlisting}

\subsection{Vectorized / Cache Optimized Approach}
The naive multiplication accesses Matrix B column-wise, which causes frequent CPU cache misses. By transposing B, we access it row-wise (contiguous memory), simulating vector operations.

\begin{lstlisting}[language=Java, caption={Vectorized (Transposed) Implementation}]
// Transpose B first
double[][] B_transposed = new double[n][n];
for(int i=0; i<n; i++) {
    for(int j=0; j<n; j++) {
        B_transposed[i][j] = B[j][i];
    }
}

// Multiply using dot product of two rows (Cache Friendly)
for (int i = 0; i < n; i++) {
    for (int j = 0; j < n; j++) {
        C[i][j] = vectorizedDotProduct(A[i], B_transposed[j]);
    }
}
\end{lstlisting}

% ------------------------------------------------------------
\section{Results and Discussion}

The benchmarks were run for matrix sizes $N \in \{256, 512, 1024, 2048, 4096\}$.

\subsection{Quantitative Analysis}

Table \ref{tab:results} shows the execution times for the largest dataset ($N=4096$).

\begin{table}[H]
\centering
\caption{Performance Metrics at $N=4096$}
\label{tab:results}
\begin{tabular}{lrrr}
\toprule
\textbf{Method} & \textbf{Time (s)} & \textbf{Speedup} & \textbf{Notes} \\
\midrule
Sequential & 2034.37 & 1.0x & Baseline \\
ParallelStream & 220.15 & 9.24x & Best Parallel \\
Threads (Manual) & 216.08 & 9.41x & Overhead overhead \\
Atomic BagTasks & 206.71 & 9.84x & Balanced Load \\
\textbf{Vectorized (Sim)} & \textbf{52.30} & \textbf{38.89x} & \textbf{Best Overall} \\
\bottomrule
\end{tabular}
\end{table}

\subsection{Visual Analysis}

\begin{figure}[H]
    \centering
    \includegraphics[width=0.95\textwidth]{task3_results.png}
    \caption{Execution Time and Speedup vs Matrix Size. Note the massive performance gain of the Vectorized approach.}
    \label{fig:results}
\end{figure}

\subsection{Discussion}

\subsubsection{Parallelism Efficiency}
The parallel implementations achieved a Speedup of approx $\sim 9.5x$. 
$$Efficiency = \frac{Speedup}{Cores}$$
Assuming a 4-core/8-thread environment, an efficiency $>100\%$ (Speedup $> 4x$) suggests that the sequential version was severely hindered by other factors or that the parallel garbage collection helped. However, the result is consistent across all parallel methods.

\subsubsection{The Victory of Memory Optimization}
The most striking result is the \textbf{Vectorized version (52.3s)}, which is $\sim 4x$ faster than the parallel versions and $\sim 38x$ faster than sequential.
\begin{itemize}
    \item \textbf{Why?} Matrix multiplication is \textit{memory-bound}. In the naive approach, fetching \texttt{B[k][j]} jumps across memory, causing Cache Misses.
    \item By transposing B, we read contiguous blocks of memory. The CPU prefetcher works effectively, keeping the pipeline full.
    \item This demonstrates that \textbf{optimizing data access patterns} often yields higher gains than simply adding more threads.
\end{itemize}

% ------------------------------------------------------------
\section{Conclusions}

Task 3 demonstrated the capabilities of Java's concurrency framework. While \texttt{ParallelStreams} offered the best balance between code simplicity and performance for parallel execution, the experiment proved that software optimization (specifically memory layout) is paramount. The single-threaded Vectorized implementation outperformed the multi-threaded naive implementations by a factor of 4, highlighting that \textbf{hardware parallelism cannot fully compensate for inefficient memory access}.

\end{document}